\textit{Silence} did not creak the way other ships creaked.

Euron Greyjoy had built her that way on purpose, years before any of this — every joint pegged and greased rather than nailed, every deck plank fitted so close a man could walk her length in a gale and hear nothing beneath his own feet but the sea itself. It was a small vanity, in the scheme of the vanities he permitted himself, but he had always believed a man’s ship should be as quiet as his crew, and his crew, tongueless every one of them, could offer no argument to that even if any of them had thought to disagree.

He stood at her prow now with the Blackwater still two days distant by his navigator’s reckoning, watching the Golden Company’s transports keep their laborious formation off his stern quarter — heavier ships, louder ships, ships full of ten thousand sellswords who talked entirely too much for Euron Greyjoy’s taste, and whom he intended, all the same, to deliver to Cersei Lannister’s gates precisely as contracted, because a bargain kept was worth infinitely more to him than a bargain merely made. Anyone could make a bargain. It took a particular kind of man to be trusted to keep one, and Euron had decided, some years back, that being that kind of man — when it suited him, and exactly as far as it suited him — was among the more useful reputations a man could cultivate.

He was thinking this, more or less, when Ser Roland Storm came to find him, as the man had taken to doing most evenings since the fleet turned north past Tarth — not from any love of Euron’s company, which the captain made no effort to disguise not feeling, but from the particular discipline of a man who believed an officer owed his commander a report whether or not he enjoyed delivering it.

“The Company’s grumbling again,” Roland said, without preamble, because he had also decided some time ago that preambles were wasted breath on a man like Euron Greyjoy. He was young for a captain, dark-haired, a bastard’s surname worn the way some bastards wore it — not as shame, but as a fact settled long ago and no longer worth anyone’s comment. He had earned his captaincy the hard way, in the Disputed Lands, promoted twice over better-born men because the Golden Company had never much cared whose seed a man carried so long as he held his line. “Third company mostly. They didn’t sign to fight a dead army. They signed to fight Northern levies and maybe a dragon queen’s rabble, and they’re hearing different rumors now, coming up from the ports.”

“Rumors.”

“That the Wall’s fallen. That there’s no Wall left to fall behind, once we’re past whatever’s coming down from the North. Sellswords don’t mind dying for coin, Lord Greyjoy. Most of them have made a kind of peace with dying for coin a long time ago. What they mind is dying for a cause nobody troubled to explain to them first.”

“Tell them,” Euron said, watching a gull work the wind off \textit{Silence}’s wake, “that they are not fighting a dead army. They are fighting whatever fool army stands between us and King’s Landing’s gold, same as always, and that the dead — if the dead trouble to exist at all outside a fisherwife’s fireside tale — are a problem for whoever loses the war first. Golden Company men have never yet been paid to worry about problems that aren’t theirs.”

“And if it’s true?”

“Then I expect the men who signed for coin will discover, same as the rest of us, that coin buys a great many things but never once bought a man out of dying, which every soldier who’s lived past his first battle already knows and pretends not to.” Euron turned from the rail at last, and gave Roland Storm the particular smile he reserved for men whose loyalty he had not yet decided whether to trust or to spend. “You didn’t climb up from a bastard’s birth to a captaincy by being a man who frightens easily, Ser Roland. I’d not have thought rumors from a fisherwife’s fire would trouble you either.”

“They don’t trouble me,” Roland said, which was not, Euron noted, quite the same as saying they didn’t trouble him at all. “It’s the men under me I’m minding. A captain who doesn’t mind his men’s fear is a captain who loses them the first hard night, whatever coin’s owed them.” He hesitated, the barest fraction of a beat, and Euron — who had made a long study of exactly such fractions in other men, and trusted them more than he trusted whole speeches — filed it away. “There’s also the matter of Prince Aegon, my lord. The men take comfort from him. A true Targaryen prince, riding with us, gives the whole venture a shape that ‘contracted for coin’ doesn’t.”

“Does he.” Euron let that sit a moment, savoring it, the way he savored most things that other men treated too gravely to enjoy. “Tell me, Ser Roland — between the two of us, captain to captain — do you believe the boy is who Connington swears he is?”

Roland Storm was quiet long enough that the answer became, in its own way, an answer. “I believe,” he said at last, carefully, “that it scarcely matters what I believe. The men believe it, or want to enough not to ask too hard, and a company marches better under a banner it believes in than a banner it’s merely paid to carry.”

“Spoken like a captain,” Euron said, “and not at all like a man who’s answered my question.” He was already turning back to the rail, already losing whatever small interest he’d had in Ser Roland Storm’s opinion of pretenders, real princes, and the difference — if any existed — between the two. He had larger uses for a Targaryen shape marching under his banner than the truth of the blood inside it, and larger concerns still, waiting for him below decks, that no captain of the Golden Company, however useful, was going to be permitted anywhere near.

Roland left him at the rail without further ceremony, glad enough to be dismissed, and made his way aft along the transport that carried the Third Company rather than back to his own ship’s boat waiting alongside — because there was one more report he owed himself before the night was done, even if he owed it to no one else.

He found the young man they called Griff exactly where he usually found him at this hour, standing alone near the stern lantern with a sellsword’s cloak thrown over shoulders too fine-boned to have earned it in any camp Roland had ever served in, watching the dark water the way a man watches something he is trying to convince himself is real. He was handsome in the particular, unplaceable way that made every soldier in the Company want to believe the story Jon Connington told about him — silver-touched hair kept close-cropped and dulled with dye that never quite held against salt air, violet eyes that could have come from a Volantene bolt-hole full of look-alikes or from the blood of Aegon the Conqueror himself, and Roland Storm, who had served under three separate lords who all swore certainty about things that turned out, in the end, to matter only exactly as much as an army’s willingness to die for them, had long since stopped trying to settle the question in his own mind.

“Ser Roland.” The boy — the prince, if the men were owed the courtesy, and Roland found he generally gave it to him whether or not he’d settled his own doubts — inclined his head with the careful, schooled grace of someone who had been drilled in it since childhood by a man who wanted very badly for it to be believed. “You’ve been at the Crow’s Eye’s ear again.”

“Someone has to be. He doesn’t keep company well.” Roland came to stand beside him at the rail, close enough for a private word, far enough that the watch wouldn’t trouble itself listening. “The men are asking after you more than they’re asking after the coin, these last nights. Thought you’d want to know that, whatever else is true.”

“Is that meant to comfort me?”

“Meant to warn you. A company that fights for a man fights harder than one that fights for coin, until the day the man disappoints them, and then it turns on him twice as hard for having believed.” Roland studied the water rather than the boy’s face, out of the same instinct for courtesy that had kept him from ever asking the question outright. “Whoever you are, Griff — Blackfyre, Targaryen, or some maester’s clever invention Connington’s spent a lifetime polishing — the men need you to be real enough to die for. I’d think on what that costs you, before we make land.”

The boy said nothing for a long while, and the silence, Roland thought, told him rather more than an answer would have.

\sceneBreak

He went down alone, once the last light had gone fully out of the sky, to the locked cabin behind his own, where the thing he had carried out of the Grey Waste of Sothoryos lived wrapped in oiled leather and old chains, and where he had gone, most nights of this voyage, to sit with it the way another man might sit with a lover he wasn’t yet certain he trusted enough to touch.

Dragonbinder, the old tales called it, though the men who’d sold it to him — dying men, all of them, which he had thought at the time an excellent recommendation and had come, since, to reconsider — had called it other names too, in tongues he did not speak and had not troubled to have translated, on the theory that a curse understood too well sometimes refused to be picked up at all.

He unwrapped it slowly, the way he always did, letting the leather fall back from horn the color of old blood, banded in red gold and black iron, graven with valyrian glyphs that seemed, in poor light, to shift very slightly whenever he was not looking at them directly. He did not sound it. He had not yet dared to sound it, not truly, not with his whole chest behind the breath the tales insisted it demanded — only once, softly, alone, testing the weight of it against his lips the way a man tests a blade’s edge against his thumb before he trusts it to a fight, and even that testing had cost him three days of a ringing in his skull that no maester’s draught had touched, and a dream, recurring since, of something vast and winged turning slowly in the dark to look for whoever had dared to call it.

He was not a man given to fear. He had told himself that so many times across so many years that he had long since stopped being able to tell whether it remained true or had simply calcified into a habit indistinguishable from truth. But he sat with the horn some while longer that night, turning it in his hands, thinking of dragons he had not yet seen with his own eyes and meant, before this war finished with all of them, to bend to a will that was not their rider’s — and for once, in the dark, alone, with ten thousand sellswords and a queen’s gold on the far side of two days’ sailing, Euron Greyjoy let himself consider, seriously and at length, exactly what a thing might cost a man before it agreed to be mastered.

He wrapped it again before he let the thought go any further, and went up on deck to stand once more at \textit{Silence}’s silent prow, and did not look back down at the locked door behind him, though he was aware, the rest of that watch, of exactly how far behind him it was.